\chapter{Фазовые диаграммы}\label{ch:diagrams}

Фазовая диаграмма --- карта: при данных температуре и составе она говорит, какие фазы устойчивы, каков их состав и
сколько каждой. Глава~\ref{ch:solutions} показала, как одна граница --- линия растворимости --- получается из кривых
$G(x)$ общей касательной. Здесь та же идея даст целые диаграммы, а с ними --- язык, на котором написаны все
справочники по сплавам.

\section{Правило фаз}

\textbf{Фаза} --- однородная часть системы с определённой структурой и составом: жидкость, $\alpha$-Zr, $\delta$-гидрид.
\textbf{Компоненты} --- независимые химические составляющие (Zr и H; Cu и Ni). Сколько параметров можно менять,
не нарушая равновесия фаз, говорит \textbf{правило фаз Гиббса}:
\begin{keyf}[Правило фаз]
\[
F = C - P + 2,\qquad \text{при постоянном давлении}\quad F = C - P + 1,
\]
$C$ --- число компонентов, $P$ --- число сосуществующих фаз, $F$ --- число степеней свободы.
\end{keyf}
Откуда оно? Каждая фаза описывается температурой, давлением и $C - 1$ независимыми концентрациями: всего
$P(C-1) + 2$ переменных. Равновесие требует равенства химического потенциала каждого компонента во всех фазах ---
$C(P-1)$ уравнений. Разность и есть $F$.

Для чистого металла при постоянном давлении ($C = 1$) две фазы сосуществуют при $F = 0$ --- только в одной точке по
температуре: плавление и полиморфные превращения идут при постоянной $T$. В двойной системе ($C = 2$) две фазы
сосуществуют в интервале температур ($F = 1$), а три --- лишь в точке ($F = 0$): эвтектика, эвтектоид, перитектика.

\section{Чистый металл}

Для одного компонента диаграмма строится в координатах давление--температура. При каждой $T$ устойчива фаза с
наименьшей $G$; линия равновесия двух фаз --- там, где их $G$ равны (рис.~\ref{fig:d1_unary}а). Наклон линии на
плоскости $(p,T)$ задаёт уравнение Клапейрона--Клаузиуса $\dd p/\dd T = \Delta H/(T\Delta V)$: если высокотемпературная
фаза плотнее, линия наклонена влево --- давление понижает температуру превращения.

\fig{d1_unary}{а) Энергии Гиббса $\beta$-Zr и жидкости относительно $\alpha$-Zr (теплоты превращения около 4 и
21~кДж/моль считаются постоянными): пересечения дают 863 и 1855\degC; серая полоса --- устойчивая фаза.
б) Схема диаграммы давление--температура циркония: $\alpha$, $\beta$, жидкость и фаза высокого давления $\omega$;
положения линий при высоких давлениях условны.}

У циркония при атмосферном давлении две твёрдые фазы: ГПУ $\alpha$ до 863\degC{} и ОЦК $\beta$ выше. ОЦК фаза
устойчива при высокой температуре благодаря энтропии: её решётка «рыхлее» для колебаний, у неё больше колебательная
энтропия, и член $-TS$ в конце концов перевешивает.

\section{Изоморфная система и правило рычага}

Если два металла неограниченно растворимы друг в друге и в жидком, и в твёрдом состоянии, диаграмма --- простейшая,
\textbf{изоморфная}. Классический пример --- медь--никель; среди систем циркония так ведут себя Zr--Hf, Zr--Ti и
(в $\beta$-области) Zr--Nb. На рис.~\ref{fig:d2_isomorph} диаграмма Cu--Ni посчитана по модели идеальных растворов
(глава~\ref{ch:solutions}) только из температур и теплот плавления чистых металлов --- и выглядит так же, как
экспериментальная.

\fig{d2_isomorph}{а) Диаграмма Cu--Ni в модели идеальных растворов: ликвидус и солидус, между ними --- двухфазная
область $L+\alpha$. Конода при 1250\degC{} для сплава с 43\,\% Ni и правило рычага. б) Кривые $G(x)$ жидкости и твёрдого
раствора при 1250\degC{} и общая касательная: точки касания --- составы $x_L$ и $x_S$ на концах коноды.}

Над \textbf{ликвидусом} --- только жидкость, под \textbf{солидусом} --- только твёрдый раствор, между ними обе фазы. В
двухфазной области составы фаз определены однозначно: проведём при данной $T$ горизонталь --- \textbf{коноду} --- до
пересечения с ликвидусом и солидусом. Это и есть точки общей касательной (рис.~\ref{fig:d2_isomorph}б). Количество
фаз даёт баланс массы: если сплав состава $x_0$ состоит из доли $f_L$ жидкости состава $x_L$ и доли $1-f_L$ твёрдого
раствора состава $x_S$, то $x_0 = f_Lx_L + (1-f_L)x_S$, откуда
\begin{keyf}[Правило рычага]
\[
f_L = \frac{x_S - x_0}{x_S - x_L},\qquad f_S = \frac{x_0 - x_L}{x_S - x_L}.
\]
\end{keyf}
Доля фазы пропорциональна длине \emph{противоположного} плеча коноды --- как у рычага с грузами. Для сплава с
43\,\% Ni при 1250\degC{} $x_L\approx0.40$, $x_S\approx0.47$, и жидкости 56\,\%.

\begin{idea}
При охлаждении сплава в двухфазной области и жидкость, и твёрдое непрерывно меняют состав: первые кристаллы
богаче тугоплавким компонентом, последние --- беднее. В равновесии диффузия выравнивает состав твёрдого. При
быстром затвердевании диффузия в твёрдом не успевает, и в зёрнах остаётся \textbf{ликвация} --- центр
богаче никелем, края беднее. Тот же сюжет повторяется с водородом в цирконии: при быстром охлаждении он не
успевает дойти до гидридов, и раствор остаётся пересыщенным.
\end{idea}

\textbf{Правила Юм-Розери} подсказывают, когда растворимость в твёрдом состоянии неограниченна: одинаковая
кристаллическая решётка, разница атомных радиусов не больше 15\,\%, близкая электроотрицательность, одинаковая
валентность. Zr и Hf (радиусы 0.160 и 0.159~нм, обе ГПУ и обе ОЦК) удовлетворяют всем правилам. Ниобий
(0.146~нм, ОЦК) неограниченно растворяется в $\beta$-Zr, но в ГПУ $\alpha$-Zr --- лишь до долей процента.

\section{Эвтектика}

Пусть в твёрдом состоянии компоненты растворяются друг в друге ограниченно --- кривая $G(x)$ твёрдого раствора
имеет две ямы (расслоение, глава~\ref{ch:solutions}). Тогда при каждой температуре нижняя выпуклая оболочка
кривых $G_L$ и $G_S$ может состоять из нескольких кусков: однофазных участков (оболочка совпадает с кривой) и
двухфазных (прямые касательные). Проделаем это при каждой температуре --- и диаграмма построена сама
(рис.~\ref{fig:d3_eutectic}).

\fig{d3_eutectic}{Эвтектическая система A--B, рассчитанная из кривых $G(x)$: жидкость --- идеальный раствор, твёрдые
$\alpha$ и $\beta$ --- одна кривая регулярного раствора с сильным расслоением ($\Omega = 16$~кДж/моль). а) Кривые $G$
чуть ниже эвтектической температуры: нижняя выпуклая оболочка касается двух ям $G_S$, жидкость выше неё. б) Диаграмма,
полученная выпуклой оболочкой при каждой температуре: эвтектика E при 452\degC{} и 58\,\% B.}

В точке E три фазы --- жидкость, $\alpha$ и $\beta$ --- лежат на одной касательной. Это нонвариантное равновесие
($F = 0$): при охлаждении жидкость эвтектического состава при одной температуре превращается в смесь двух твёрдых
фаз, $L\to\alpha+\beta$. Обычно смесь получается тонкими чередующимися пластинками (\textbf{эвтектика}): атомам
выгоднее перераспределяться на короткие расстояния между соседними пластинками.

Другие нонвариантные превращения двойных систем:
\begin{center}\small
\begin{tabularx}{\linewidth}{>{\raggedright\arraybackslash}p{3.1cm}>{\raggedright\arraybackslash}p{3.6cm}X}
\toprule
превращение & реакция при охлаждении & пример \\
\midrule
эвтектическое & $L\to\alpha+\beta$ & Al--Si, Pb--Sn, Zr--Fe \\
эвтектоидное & $\gamma\to\alpha+\beta$ (всё твёрдое) & Fe--C (727\degC), \textbf{Zr--H: $\beta\to\alpha+\delta$ при $\approx550$\degC} \\
перитектическое & $L+\alpha\to\beta$ & Fe--C при 1495\degC, Cu--Zn \\
монотектоидное & $\beta_1\to\alpha+\beta_2$ & \textbf{Zr--Nb: $\beta_{\text{Zr}}\to\alpha_{\text{Zr}}+\beta_{\text{Nb}}$ около 620\degC} \\
перитектоидное & $\alpha+\beta\to\gamma$ & образование ряда интерметаллидов \\
\bottomrule
\end{tabularx}
\end{center}

\section{Кривые охлаждения и структура}

Как проверить диаграмму? Самый старый способ --- \textbf{термический анализ}: охлаждать сплав и записывать
температуру. Пока превращения нет, температура падает плавно. Превращение выделяет теплоту и задерживает
охлаждение (рис.~\ref{fig:d4_cooling}а):
\begin{itemize}
\item чистый металл и эвтектический сплав дают \textbf{площадку} --- нонвариантное превращение идёт при постоянной
  температуре;
\item сплав между чистым компонентом и эвтектикой даёт \textbf{излом} на ликвидусе (выделение первичных кристаллов
  растягивается на интервал температур) и площадку при $T_E$, где кристаллизуется оставшаяся жидкость.
\end{itemize}

\fig{d4_cooling}{а) Расчётные кривые охлаждения (теплоотдача пропорциональна $T - T_{\text{среды}}$, выделение теплоты
кристаллизации --- по долям фаз из диаграммы рис.~\ref{fig:d3_eutectic}). б) Доли фаз при охлаждении сплава с 20\,\% B по
правилу рычага: первичная $\alpha$ растёт от ликвидуса, при $T_E$ оставшаяся жидкость (18\,\%) превращается в
эвтектику. в) Схема структуры этого сплава: зёрна первичной $\alpha$ в пластинчатой эвтектике $\alpha+\beta$.}

Правило рычага сразу говорит, какой будет структура. В сплаве с 20\,\% B непосредственно над $T_E$ жидкость
имеет эвтектический состав 58\,\%, первичная $\alpha$ --- 11\,\%; жидкости $(20-11)/(58-11)\approx18\,\%$. Значит,
в готовом сплаве 82\,\% первичных зёрен и 18\,\% эвтектики (рис.~\ref{fig:d4_cooling}в). Тот же расчёт, сделанный по
снимку шлифа «наоборот», позволяет оценить состав.

\section{Твёрдые растворы и соединения в сплавах циркония}

Твёрдые растворы бывают двух видов. В \textbf{растворе замещения} атомы второго компонента занимают узлы решётки (Sn,
Nb, Fe, Hf в цирконии). В \textbf{растворе внедрения} маленькие атомы сидят в пустотах (H, O, N, C). Растворимость
внедрения обычно мала, но водород исключение по подвижности, а кислород --- по растворимости: в $\alpha$-Zr его можно
растворить до 30~ат.\,\%, и он сильно упрочняет металл.

Легирующие элементы делятся по тому, как они сдвигают превращение $\alpha\leftrightarrow\beta$:
\begin{itemize}
\item \textbf{$\alpha$-стабилизаторы} (O, Sn, Al) расширяют область $\alpha$ вверх --- граница $\alpha/(\alpha+\beta)$ идёт
  к более высоким температурам;
\item \textbf{$\beta$-стабилизаторы} (Nb, Fe, Cr, Ni, Mo, H) опускают её; часто с эвтектоидом или монотектоидом. Водород ---
  сильный $\beta$-стабилизатор: в системе Zr--H область $\beta$ опускается до эвтектоида при $\approx550$\degC.
\end{itemize}
Когда растворимость превышена, выделяются \textbf{интерметаллиды} --- соединения с собственной решёткой и узким составом.
В сплавах циркония это, например, фазы Лавеса Zr(Fe,Cr)$_2$ в циркалое и Zr(Nb,Fe)$_2$ в ниобийсодержащих
сплавах (в том числе Э635), а также частицы $\beta$-Nb. Мелкие (десятые доли микрометра) частицы влияют на
коррозию, на рекристаллизацию и могут служить местами зарождения гидридов.

\section{Выделение из пересыщенного раствора}

Пусть растворимость второго компонента падает с температурой (наклонный \textbf{сольвус}). Сплав нагревают в
однофазную область (\textbf{закалка на твёрдый раствор}), быстро охлаждают --- раствор остаётся пересыщенным, --- а
затем выдерживают при умеренной температуре (\textbf{старение}). Избыток выделяется в виде частиц второй фазы.
Пересыщение даёт движущую силу, а скорость задают зарождение и диффузия (главы~\ref{ch:nucleation} и
\ref{ch:diffusion}). Часто сначала выделяются не равновесные, а метастабильные фазы с меньшим барьером
зарождения: когерентные зоны, переходные фазы, и лишь потом --- равновесная.

Водород в цирконии ведёт себя точно так же, только всё происходит быстрее: подвижность водорода огромна, и
«старение» идёт прямо при охлаждении. Линия TSSD --- сольвус; пересыщение копится при охлаждении ниже неё; выпадение
начинается на TSSP, где пересыщение достаточно для зарождения. Метастабильный $\gamma$-гидрид при быстром охлаждении ---
аналог переходных фаз.

\section{Диаграмма Zr--H со стороны циркония}

Полная диаграмма Zr--H сложна: несколько гидридов ($\gamma$-ZrH, $\delta$-ZrH$_{\sim1.6}$, $\varepsilon$-ZrH$_2$),
высокотемпературная $\beta$-фаза, эвтектоид. Для оболочек важен только угол со стороны циркония
(рис.~\ref{fig:f2_zrh}):
\begin{itemize}
\item $\alpha$ --- твёрдый раствор водорода в ГПУ цирконии;
\item $\alpha+\delta$ --- раствор плюс гидрид; граница --- линия растворимости (TSSD);
\item $\beta$ --- ОЦК цирконий; водород её стабилизирует, поэтому область $\alpha+\beta$ клином уходит вниз;
\item эвтектоид при $\approx 550$\degC: $\beta \to \alpha + \delta$. Растворимость в $\alpha$ там максимальна,
  $\approx 6$~ат.\,\% (около 690~ppm).
\end{itemize}

\begin{figure}[htbp]\centering
\includegraphics[width=0.6\linewidth]{f2_zrh}
\caption{Угол диаграммы Zr--H со стороны циркония (схема; сольвус $\alpha/(\alpha+\delta)$ --- по линии TSSD,
остальные границы условно прямыми). Вертикальные пунктиры --- 100, 300 и 600~ppm водорода.}
\label{fig:f2_zrh}
\end{figure}

\begin{warnbox}
Сплавы (Nb, Sn, Fe, O) немного сдвигают линии, а линия выпадения зависит от режима охлаждения. Поэтому для
конкретного материала линии берут из его собственных измерений, а не из общей диаграммы. И помните: TSSP --- не
линия равновесной диаграммы, а кинетическая граница, о которой равновесная термодинамика ничего не знает
(главы~\ref{ch:solutions} и \ref{ch:nucleation}).
\end{warnbox}

\section{Сколько выпадает гидрида}

По правилу рычага при комнатной температуре почти весь водород сидит в гидриде (растворимость там --- доли ppm).
В $\delta$-ZrH$_{1.66}$ водорода $\approx 1.8$\,\% по массе (18\,000~ppm). Значит, при 160~ppm массовая доля
гидрида $160/18\,000 \approx 0.9\,\%$, а объёмная --- около 1\,\% (гидрид менее плотен). На травленом шлифе гидриды
кажутся занимающими в несколько раз большую долю площади --- травление расширяет их следы (глава~\ref{ch:methods}).
Об этом стоит помнить, сравнивая снимки с расчётом.

\begin{exercise}
Сколько степеней свободы у системы Zr--H в точке эвтектоида при постоянном давлении? А в области $\alpha+\delta$?
Что это значит для температуры, при которой начинается выпадение гидрида в сплаве с данным содержанием водорода?
\end{exercise}

\begin{exercise}
По данным рис.~\ref{fig:d2_isomorph} ($x_L = 0.40$, $x_S = 0.47$ при 1250\degC) найдите долю жидкости в сплаве с 45\,\% Ni.
\end{exercise}

\begin{exercise}
Найдите долю эвтектики в сплаве с 30\,\% B на рис.~\ref{fig:d3_eutectic} ($x_\alpha = 0.11$, $x_E = 0.58$ при $T_E$).
\end{exercise}

\begin{exercise}
Сколько гидрида (по массе) выпадет при охлаждении от 400\degC{} до комнатной температуры в Э635 с 300~ppm водорода?
Какая доля водорода выпадет до 300\degC, если выпадение идёт по TSSD (медленное охлаждение)?
\end{exercise}
